\section{Applications}
\label{sec:applications}

\review{None of the examples use variants, which is arguably the most
important feature that \dargent provides which surpasses the expressiveness
of C.}

\ambroise{we do use variants (to account for enums), although with empty constructor payloads. If I
  have time to implement it, the
  maybe example suggested by the reviewer could be more convincing.
}

In this section, we showcase how \Dargent helps in developing and formally
verifying systems code via some small examples, the full source code of which is available in the supplementary material.
 
% \review{more information about effort involved (especially something quantitative, if possible), and possibly information about performance. also for \S~6.}


\subsection{Power Control System}\label{ssec:power-control-system}

To demonstrate the improved readability of programs that \Dargent offers, 
we reimplemented the power control 
system for the STM32G4 series of ARM Cortex
microcontrollers by ST Microelectronics, based on a C implementation
from LibOpenCM3\footnote{\url{http://libopencm3.org/}}, an open-source low-level
hardware library for ARM Cortex-M3 microcontrollers.
The original C file\footnote{\url{http://libopencm3.org/docs/latest/stm32g4/html/pwr_8c_source.html}}
is about one hundred lines long, and consists of eight functions of a few lines
each that perform some bitwise operations on the device register. 
% \review{how long is the corresponding \cogent code?}
The \cogent implementation is a bit smaller as each function is now one line long. The
most involved part is the \dargent layout itself, which is derived from the hardware specifications of the device.

The 32-bit power control register holds several pieces of information, each occupying a certain
number of bits in the register. We define the register as a record type in \cogent
and use \dargent to prescribe the placement of each field. With a \dargent layout,
the functions in the power control code no longer involve bitwise operations, but merely set values to individual fields.
For example, the function to set the voltage output scaling bit in C:
\begin{lstlisting}[language=C]
  void pwr_set_vos_scale(enum pwr_vos_scale scale) {
    uint32_t reg32;
  
    reg32 = PWR_CR1 & ~(PWR_CR1_VOS_MASK << PWR_CR1_VOS_SHIFT);
    reg32 |= (scale & PWR_CR1_VOS_MASK) << PWR_CR1_VOS_SHIFT;
    PWR_CR1 = reg32;
  }
\end{lstlisting}
translates to the \cogent function setting the $\FieldN{vos}$ field in the record:
\begin{lstlisting}[language=Cogent]
  pwr_set_vos_scale : (Cr1, Vos_scale) -> Cr1
  pwr_set_vos_scale (reg, scale) = reg { vos = scale } 
\end{lstlisting}
It can be seen from this example that \cogent allows systems programmers to write code
on an abstract level, while retaining low-level control of the implementation details with \dargent{}.


% for the STM32G4 series of ARM Cortex Microcontrollers by ST Microelectronics.

% \review{Furthermore, while Dargent provides a lot of control over data *layout*, it does
% not (seem to) provide control over the exact *access pattern* that is used to
% get and set fields. Of course, usually that is an implementation detail, but
% several of the examples involve using Dargent to describe the layout of hardware
% registers (which I assume refers to hardware that is controlled via MMIO).
% Hardware register access in C is usually done with `volatile`, ensuring that the
% access happens exactly as written (e.g., a single 32-bit write, even if only the
% first 16 bits were changed). For example: presumably, in pwr\_set\_vos\_scale,
% PWR\_CR1 is a volatile 32-bit integer? Does the code generated by Dargent
% guarantee that the setter will do a single 32-bit volatile write? The
% description of the generated setters and getters is rather vague, so it might
% well be the case that Dargent generates a 16-bit write when only some bit in the
% first half of PWR\_CR1 changed, which might or might not actually be supported be
% the hardware this driver is controlling. Or, worse, it could generate multiple
% small writes. In these settings where the exact sequence of memory writes and
% their sizes matters, how does a Dargent user ensure that the generated code does
% what it is supposed to do?}

% \zilinc{Do we talk about it here, or early in \S~3?}
% 
% \ambroise{I think we should sketch the access pattern in sec 3.4 (and discuss
%   performance, in particular regarded accessing nested subfields) on custom
%   getters and setters, but discuss the volatile issue here.}


%\paragraph{Dealing with volatile registers}
%This example also demonstrates a current limitation of our system:
%the verification framework does
%not provide any guarantees about the exact access pattern of
%getters and setters. 
%This is however crucial when dealing with hardware registers, since
%the write (or read) order can matter.
% Consider, for example, the following
%\cogent function from the power control system.
%\begin{lstlisting}[language=Cogent]
%pwr_enable_power_voltage_detect : (Cr2, Pvd_level) -> Cr2
%pwr_enable_power_voltage_detect (reg, pls) = reg { pls = pls, pvde = True }
%\end{lstlisting}
%This code actually compiles to two subsequent writes to the register CR2:
%it first sets the programmable voltage detector level \lstinline|pls|
%and then enables the programmable voltage detector (by setting the boolean
%flag).
%If the writes were written in reverse order, however, the detector would run
%with an unknown selection level between the two writes.
%% Although the device specification does not explicitly
%% discourage such a write pattern,
%To avoid this kind of unwanted behaviour, sequentialisation is
%usually avoided when dealing with hardware registers. 
%In fact, the original C implementation performs a single write to the register:
%\begin{lstlisting}[language=C]
%void pwr_enable_power_voltage_detect(uint32_t pvd_level)
%{
%  uint32_t reg32;
%  reg32 = PWR_CR2;
%  reg32 &= ~(PWR_CR2_PLS_MASK << PWR_CR2_PLS_SHIFT);
%  PWR_CR2 = (reg32 | (pvd_level << PWR_CR2_PLS_SHIFT) | PWR_CR2_PVDE);
%}
%\end{lstlisting}
%%In future work, we plan to extend \cogent so that it can perform
%%multiple field updates/reads in a single operation, when possible.

% The difficulty
% Beyond the implementation, a  difficulty relates to
% the formalisation of the different write patterns, which are
% currently indiscernable because of the memory model of the
% AutoCorres library on which \cogent's verification framework
% is based.

% As an example, flipping multiple booleans flags
% will always result in multiple writes
% However, hardware register access typically requires 
% that  depends C is usually done with `volatile, ensuring that the
% access happens exactly as written (e.g., a single 32-bit write, even if only the
% first 16 bits were changed).



\subsection{Bit Fields}\label{ssec:bitfields}

In systems code, there is a regular pattern consisting of declaring a structure whose fields have non-standard
 bit widths, as in the following example taken from a CAN driver%
 \footnote{\url{https://github.com/seL4/camkes-vm-examples/blob/89f5d7b7ac373c8e9f000e80b91611e561358ef6/apps/Arm/odroid_vm/include/can_inf.h\#L34}}
 where the first field is an identifier on 29 bits.
\begin{lstlisting}[language=C]
  struct can_id {
    uint32_t id:29;
    uint32_t exide:1;
    uint32_t rtr:1;
    uint32_t err:1;
  };
\end{lstlisting}
Although this is not supported by the AutoCorres library~\citep{autocorres} on which the \cogent
verification framework is based, we can simulate this feature using \Dargent.

% As a first attempt, consider the following \cogent type.
% \begin{lstlisting}
%   type R = { f1 : Bool, f2 : U8 }
%      layout record { f1 : 1b, f2 : 4b }
% \end{lstlisting}
% The \cogent compiler would actually complain that a \lstinline{U8} value does not fit in 4 bits.
% % As far as the verification is concerned,
% This mismatch compromises the first get/set
% lemma (see~\S\ref{sec:get-set-lemmas}), since setting the field \lstinline{f2}
% with a \lstinline{U8} value $v$ and then retrieving
% it does not yield $v$, but its truncation on 4 bits.
% To solve this issue, we need to introduce a new type \lstinline{U4} in \cogent. 
% This type essentially translates to \lstinline{U8} in C, but the value relation 
% for this type ensures that the value fits on 4 bits.
The last three fields are one bit long and can thus be handled with a \cogent
boolean type.
In order to cover the 29-bit field, we can use a primitive 29-bit integer $\PrimType{U29}$. These non-word-size integers are a new extension we added to \cogent for this purpose, and consisted only of a dozen changed lines of code. We compile such integers to
the smallest standard integer type that can contain the type, since the C language does not
natively support such types. For instance, $\PrimType{U7}$ is
compiled to $\PrimType{U8}$, while $\PrimType{U20}$ is compiled to $\PrimType{U32}$.
Thanks to the extension of the value relation (see \autoref{ssec:record-value-rel}) 
to those new types, compiled \cogent programs maintain the invariant that C
values always fit in the narrower bit-width. 
%
%
%
% However there is no built-in type
% that can be used for the 29-bit field (recall the match rules in~\S\ref{ss:formalised-semantics}). 
% Thankfully, \cogent provides a modular mechanism to introduce new types, as abstract types,
% whose definitions must be provided by the programmer in C.
% The verification framework must also be extended to take into account abstract values for these types.
% % , and provide the value relations for these new types. 
% For example, in the shallow embedding, we embed $\ConsN{U29}$
% as the Isabelle type of natural numbers that fit in 29 bits.
%
% By first declaring the abstract type $\ConsN{U29}$,
With a \dargent layout, we can define a \cogent version of the above C structure:
\begin{lstlisting}[language=Cogent]
  type CanId = { id : U29, exide : Bool, rtr : Bool, err : Bool }
    layout record { id : 29b, exide : 1b, rtr : 1b, err : 1b}
\end{lstlisting}
% Abstract types are implicitly boxed, hence the additional
% unboxed annotation \lstinline{#} in front of $\ConsN{U29}$.

% When encountering an unboxed abstract type in a layout, the compiler does not check
% any compatibility between the bit range and the type. In the generation of getters and setters,
% the field is dealt as if it were
% an ordinary integer type, with a final C cast to the abstract type.
% For some technical reason, abstracts types must always be embedded in
% a custom structure. As a consequence, we introduce an additional layer of compilation 
% that replaces this brutal C cast with a more explicit one.
% Nicely enough, our tactics are still able to prove the get/set lemmas
% stated in~\S\ref{sec:get-set-lemmas}.
% 

% To demonstrate how this approach works in practice, we proved functional correctness
% of the following simple function that updates the second or third field
% depending on the value of the first field.
% \begin{lstlisting}
%   foo : R -> R
%   foo r =
%     if r.f1 !r then
%       let v = u8_to_u4 (u4_to_u8 r.f3 .&. 0x0c) !r in
%       r {f3 = v}
%     else
%       let v = u8_to_u2 (u2_to_u8 r.f2 + 1) !r in
%       r {f2 = v}
% \end{lstlisting}
% Abstract types alone are not very useful since no \cogent operation applies to them.
% We thus extend our example with back and forth casts between $\ConsN{U29}$, and $\ConsN{U32}$,
% declared in \cogent as \emph{abstract functions}. This means that we implement them in C, 
% and we also provide their semantics at each layer of the refinement proof, and 
% prove properties about them~\citep[\S3.3.2]{OConnor:2016}.
% \begin{lstlisting}
%   u4_to_u8 : #U4 ->  U8
%   u8_to_u4 :  U8 -> #U4
%   u2_to_u8 : #U2 ->  U8
%   u8_to_u2 :  U8 -> #U2
% \end{lstlisting}
% These functions are abstract: they are implemented in C. We exploit the design
% of the verification framework that eases the extension of the semantics to account
% for abstract functions. 

% \review{how much manual effort is needed? hours? LoC? etc.}
% Although some manual verification effort is required (about 700 lines of Isabelle/HOL proof) on these irregular-sized integers and the
% cast functions, it is only a one-off effort, and they can be mechanically produced.
% We leave it as future work to have native support for these integer types from the \cogent language
% and the compiler.


\noindent Bit fields can play an important role in systems programming. This example demonstrates that 
\dargent not only allows programmers to store data in a compact manner, but also provides a 
principled way to reason about C bit fields without needing to extend our C verification tools
to support the bit field feature in C.


\subsection{Custom Getters and Setters}\label{ssec:partial-rec}

Suppose we are writing a \cogent system that must work on a complicated, externally defined C structure that involves many fields, but the \cogent component is only concerned with a few specific fields. 

Because \dargent makes all record accesses go through getter and setter functions, we can reuse this mechanism by providing \emph{custom} getter and setter functions in C, that extract (or write to) the relevant fields inside the large C structure.  Then, we can represent this large C structure as a simple \cogent record with only the relevant fields.
%% \begin{lstlisting}[language=Cogent]
%%   type Entry = { id : U32, value : U32 }  
%% \end{lstlisting}
The custom getters and setters provided by the user must compose well (see~\autoref{sec:get-set-lemmas}).

We have successfully applied this approach on a small example, where the
C structure has the same fields as the original \cogent type, extended
with an additional field. Whilst this example is contrived, this feature
does have real application: It makes it possible to \emph{inherit}
pre-defined C data structures. For example, in a previous \cogent implementation of a Linux \texttt{ext2fs} driver, \citet{Amani:2016}
had to define their own \Cogent \verb!Ext2Inode! type which corresponds to the standard C \lstinline[language=C]!struct ext2_inode! type, but did not include fields that were irrelevant to the \cogent implementation, such as  certain unsupported file modes and spinlocks. This required tedious glue code to marshal back-and-forth between the representations. 
Using this technique, the standard C \lstinline[language=C]!struct ext2_inode! can be used directly as the representation of the \verb!Ext2Inode! type, thus eliminating this glue code entirely.


%Currently, this customisation of getters and setters is accomplished rather awkwardly, requiring the programmer to assign the type a dummy layout to force the generation of getters and setters, and then replace the generated getters and setters with the desired custom functions.
%Of course, to get the full refinement proof, one needs to prove that getters and
%setters compose well (\S~\ref{sec:get-set-lemmas}).
% In future
%work, we plan to improve the compiler to make this approach more
%user-friendly.
